===== BEGIN FILE: DAMPE_FD.tex =====
\documentclass[aps,prl,showpacs,superscriptaddress,nofootinbib,twocolumn,10pt]{revtex4}
\usepackage{color,amsmath,amssymb,graphicx,latexsym,lineno,amsfonts}
\usepackage{threeparttable,multirow,txfonts,subfigure,booktabs}


\def\nat{Nature\ }
\def\aap{Astron.\ Astrophys.\ }
\def\aaps{Astron.\ Astrophys.\ Supp.\ }
\def\apj{Astrophys.\ J.\ }
\def\apjl{Astrophys.\ J.\ Lett.\ }
\def\apjs{Astrophys.\ J.\ Supp.\ }
\def\aj{Astron.\ J.\ }
\def\mnras{Mon.\ Not.\ Roy.\ Astron.\ Soc.\ }
\def\physrep{Phys.\ Rept.\ }
\def\pasa{Publ.\ Astron.\ Soc.\ Aust.\ }
\def\pasj{Publ.\ Astron.\ Soc.\ Japan\ }
\def\planss{Planet.\ Space\ Sci.\ }
\def\prc{Phys.\ Rev.\ C\ }
\def\prd{Phys.\ Rev.\ D\ }
\def\prl{Phys.\ Rev.\ Lett.\ }
\def\apss{Astrophys.\ Space\ Sci.\ }
\def\araa{Annu.\ Rev.\ Astron.\ Astrophys.\ }
\def\jcap{J.\ Cosmol.\ Astropart.\ Phys.\ }
\def\jgr{J.\ Geophys.\ Res.\ }
\def\ssr{Space\ Sci.\ Rev.\ }

\graphicspath{{fig/}}


\begin{document}

\title{Diverse properties of electron Forbush decreases revealed by the Dark Matter Particle Explorer}

\input{author20250812.tex}

\begin{abstract}
The Forbush decrease (FD) of cosmic rays is an important probe of the
interplanetary environment disturbed by solar activities. In this work,
we study the properties of 8 FDs electrons (including positrons) between
2 GeV and 20 GeV from January, 2016 to March, 2024, with the Dark Matter
Particle Explorer. The maximum decrease amplitudes of these events are
about $30\%-15\%$, and the amplitudes reduce with energy. The recovery
time of these events shows diverse behaviors of their energy-dependence.
Some of them show strong energy-dependence, while some have a nearly
constant recovery time. It has been shown that such diverse behaviors
could be related with the geometry of the disturbed regions of the
interplanetary space by coronal mass ejections (CME), represented by the
combined effect of the CME velocity, angular spread, and ejection direction.
\end{abstract}

\pacs{96.25.Qr,98.70.Sa}

\maketitle

{\it Introduction.} ---
Galactic cosmic rays (GCRs) are affected by heliospheric plasma and the
associated magnetic fields when propagating through the solar system,
resulting in flux variations at different time scales (e.g.,
\cite{2000SSRv...93...55C,2013LRSP...10....3P}).
One type of short-time variations, known as Forbush decreases (FDs)
\cite{1937PhRv...51.1108F,1937Natur.140..316H}, is the quick drop and
relatively slow recovery of the GCR fluxes over timescales of days to
weeks. FDs are caused by energetic solar activities such as coronal mass
ejections (CMEs) \cite{2010SoPh..263..223Y, 2000SSRv...93...55C} or
corotating interaction regions (CIRs) \cite{1999AdSpR..23..567H},
which disturb the interplanetary space and block part of GCRs from
reaching the detector. Fast CMEs can drive strong interplanetary shocks
that also play an important role in generating FDs, such as the two-step
FD including an initial decrease by the shock and a further reduction
when the CME itself arrives \cite{2011JGRA..11611103J}. FDs have also
been observed far away from the Earth
\cite{2018A&A...611A..79G,1993JGR....98....1B,2019ApJ...874..167A},
indicating that they are universal phenomenon in the solar system.

FDs are widely studied by ground-based neutron monitors (NMs) or muon
detectors (MDs), which measure secondary particles produced by interactions
between hadronic GCRs and the atmosphere. The NM or MD data show the
integral variations of the GCR intensities above specific energies,
which depend on locations of the detectors due to a latitude dependent
geomagnetic vertical rigidity cutoff (VRC; \cite{2005AdSpR..36.2012S}).
Nonetheless, direct measurements by space experiments, in spite that
they are limited by effective areas, provide better composition and
energy resolution, and are thus crucial to understanding the physics
of FDs. Very few direct measurements of FDs were available, by balloon
experiments and space satellites \cite{1961JGR....66.3950M,2018ApJ...853...76M,2021ApJ...920L..43A,2023SoPh..298....9L}.
The PAMELA experiment studied FDs of protons, helium nuclei, and
electrons, showing charge-sign dependences of their recovery properties
\cite{2018ApJ...853...76M}. Further investigation of a sample of proton
FDs using PAMELA data showed that there are two types of recovery times,
type I with significant energy dependence and type II without energy
dependence \cite{2023SoPh..298....9L}. Similarly, ground-based detectors
have identified two distinct types of recoveries
\cite{2007AdSpR..40..342J,2008JGRA..113.7102U,2016ApJ...827...13Z},
whose interpretation remains unclear. Using the daily proton fluxes
measured by AMS-02 \cite{2021PhRvL.127A1102A}, Wang et al. studied the
properties of FDs caused by CMEs and CIRs, finding that they do not differ
much from each other \cite{2023ApJ...950...23W} .

Observations of FDs of the leptonic component are even rarer. Daily fluxes
of electrons and positrons from May 2011 to November 2021 were published by
AMS-02 \cite{2023PhRvL.130p1001A,2023PhRvL.131o1002A}, but no detailed
analysis of the FDs has been carried out. An electron FD was reported by
PAMELA \cite{2018ApJ...853...76M}, shows a faster recovery than protons and
helium nuclei below 2 GV. The properties of an FD of cosmic ray electrons and
positrons (hereafter CREs) were measured by the Dark Matter Particle Explorer
(DAMPE) with significantly improved statistics and wide energy coverage,
showing clearly energy-dependent recovery times between 2 and 20 GeV
\cite{2021ApJ...920L..43A}. Note that the DAMPE detector itself does not
distinguish electrons from positrons due to the lack of magnetic field, and in
this work we add electrons and positrons together for the study. The fraction
of positrons in the sample is a few percents in our interested energy range
\cite{2013PhRvL.110n1102A}. The decrease and recovery behavior observed by
DAMPE for this event can be properly accounted for using a diffusion
barrier model \cite{2018ApJ...860..160L,2021ApJ...920L..43A}, the more clear
underlying physics of FDs needs further exploration using more events. It is
an intriguing question whether the recovery times of CRE FDs also experience
two types of energy dependence as thoses of GCR nuclei.

In this work, we analyze the DAMPE CRE data from January 1, 2016 to
March 31, 2024 to study the population of FDs. A total of eight large
FDs are identified, each  with maximum decrease amplitudes exceeding $30\%$ of
the average fluxes. The decrease and recovery properties of these events will
be investigated in detail.


{\it DAMPE detector.} ---
The DAMPE is a space detector for the observations of cosmic ray nuclei,
electrons and positrons, and $\gamma$-ray photons
\cite{2017APh....95....6C,2025arXiv251105409D}.
The DAMPE detector is composed of four sub-detectors: a plastic scintillator
detector (PSD; \cite{2017APh....94....1Y}) on the top for charge measurement,
a silicon-tungsten tracker-converter (STK; \cite{2016NIMPA.831..378A}) for
trajectory and charge measurement, a Bismuth-Germanium-Oxide imaging
calorimeter (BGO; \cite{2012ChPhC..36...71Z}) for energy measurement and
electron-proton separation, and a neutron detector
(NUD; \cite{2020RAA....20..153H}) for auxiliary electron-proton separation.
Serving as the key subdetector, the BGO calorimeter consists of 14 layers,
each containing 22 bars for a total depth of about $32X_0$ which enables
$\sim1.5\%$ energy resolution for $\sim10$ GeV CREs \cite{2017APh....95....6C}.
Different from observing the corona graphs, DAMPE is very suitable for
measurements of CREs with a relatively large acceptance of $\sim0.3$ m$^2$~sr
and a high energy resolution. The detector was integrated in a dedicated
satellite platform, and was launched into a sun-synchronous orbit at an
altitude of 500 km on December 17, 2015 from the Jiuquan base
\cite{2017APh....95....6C}.

\begin{figure*}
\centering
\includegraphics[width=1.0\textwidth]{Fig1_ShowFit.eps}
\caption{Time profiles of CRE fluxes in four energy bins measured by DAMPE
in November, 2021 (blue dots), compared with the DAMPE $T_0$ rate (green
curves) and OULU NM data (red curves). The errorbars of the DAMPE CRE fluxes
include both the statistical uncertainties and systematic uncertainties,
added in quadrature. The 4 panels correspond to 4 different CRE energy bins
as labelled. Arrows in the top-left panel show the two step decreases
revealed by the NM data and DAMPE $T_0$ data, and the black dashed line
shows the arrival time of the interplanetary shock.}
\label{fig:2021FD}
\end{figure*}


{\it CRE selection and flux calculation.} ---
The housekeeping system of DAMPE records incident counts of events every 4
seconds, referred to as $T_0$ counts, to monitor the condition of the data
acquisition system \cite{2019RAA....19..123Z}. The $T_0$ trigger rate is a
direct measure of the low-energy cosmic ray intensity, which varies
from dozens Hz near the equator to tens of thousands Hz at the poles.
Simulations indicate that $T_0$ triggering occurs when the energy of a particle
is higher than 100 MeV. FD candidates are selected using the DAMPE $T_0$ rate
and the OULU NM count rate, both exhibit a rapid decrease plus a relatively
slow recovery. Specifically, if the DAMPE $T_0$ rate decreases by more than
10\% and the OULU NM count rate decreases over 3\% within 3 days, the
corresponding time window is labelled as a possible FD event occurs.
Then a cross match with the SOHO LASCO CME
catalog\footnote{https://spase-metadata.org/NASA/Catalog/SOHO/LASCO/CME\_Catalog.html} is further required.
In total we find 8 candidate FDs according to this procedure. The basic
information of these 8 FD events can be found in Table~S1 of the
{\tt Supplemental Material}. For each candidate event, we select data from
10 days before to 30 days after the starting time for the follow-up analysis.
CRE candidates are selected following the method described in
Ref.~\cite{2021ApJ...920L..43A}.
Data recorded during the passage of the South Atlantic Anomaly (SAA) region
are excluded. Both events passing the high-energy trigger (HET) and low-energy
trigger (LET) \cite{2019RAA....19..123Z} conditions are used.
The reconstructed particle charge \cite{2019APh...105...31D} is set
to be $Z\le 1.7$, which rejects $\sim99\%$ of nuclei heavier than protons.
We further require that the shower maximum is not at the edge of the first
six BGO layers, and the reconstructed track is required to penetrate the
first four BGO layers. The pre-selection efficiency is shown in Fig.~S1 of
the {\tt Supplemental Material}. We can see that the efficiency reaches about
80\% for energies higher than a few GeV. At low energies the efficiency drops
quickly, due mainly to the trigger efficiency.

After these pre-selections, the main background for CREs is protons.
We use a particle identification (PID) parameter, defined as
${\rm PID} = F(E)[\log(R_{r}) \cdot \sin\theta + \log(R_{l}) \cdot \cos\theta]$
to distinguish CREs from protons, where $F(E)$ is an energy decoupling
polynomial, ${\theta}$ is a rotation angle, ${R_{r}}$ and ${R_{l}}$ are
parameters describing the radial and longitudinal extensions of the shower
in the calorimeter \cite{2021ApJ...920L..43A}. Fig.~S2 of the
{\tt Supplemental Material} illustrates the distributions of the PID
parameter for two energy bins, $3.17-3.44$ GeV and $7.35-8.00$ GeV.
CREs would typically have ${\rm PID}\leq 2$. To properly estimate the
proton background, we carry out a fit to the data with Monte Carlo (MC)
simulation templates of CREs and protons. Since CMEs affect electrons
and protons differently \cite{2018ApJ...853...76M}, the background is
estimated independently for each energy and time bin. We find that the
proton background is about 2\% to 8\% for energies from 2 to 20 GeV, with
about $15\%$ relative changes (differ case by case) during the FD period.

The CRE flux is calculated as
\begin{equation}
    I_{i,j}=\frac{N_{i,j}}{{\Delta} E_{i} A_{i} T_{i,j} \eta_{i}},
    \label{eq2}
\end{equation}
where ${N_{i,j}}$ is the number of selected CREs after the background
subtraction in the ${i}$th energy bin and ${j}$th time bin, ${T_{i,j}}$
is the corresponding exposure time, ${\Delta E_{i}}$ is the energy bin
width, ${A_{i}}$ is the geometric acceptance, and $\eta_i$ is the total
efficiency. To avoid the shielding effect from the geomagnetic field,
we require that the particle's rigidity is higher than 1.2 times
of the VRC values. Since the VRC value varies with latitude, the rigidity
threshold results in variations of the exposure time with time and energy
(see Fig.~S3 of the {\tt Supplemental Material}). On average,
the effective exposure time is approximately 80\% of the total time for
energies above 15 GeV, and decreases gradually to about 30\% at 2 GeV.

\begin{figure*}
\centering
\includegraphics[width=1.0\textwidth]{Fig2_Amp.eps}
\caption{The decrease amplitude as a function of energy, for the eight FDs
observed by DAMPE.}
\label{fig:amp}
\end{figure*}

\begin{figure*}
\centering
\includegraphics[width=1.0\textwidth]{Fig3_Rec.eps}
\caption{The recover time as a function of energy, for the eight FD observed
by DAMPE. The solid line in each panel shows the best-fit power-law result.}
\label{fig:tau}
\end{figure*}

Since we focus on the relative change of the CRE fluxes at a time scale of
about one month, the main systematic uncertainties are the efficiency
stabilities. The efficiency fluctuation is estimated to be $<0.15\%$ within
one month after correcting the temperature effect \cite{2019APh...106...18A}.
It has also been shown that there is a long-term evolution of the trigger
efficiency due to radiation damage, electronics aging, and gain decrease
of photomultipliers, which results in a decrease of the HET efficiency by
about 5\% per year at 2 GeV and smaller with increasing energies
\cite{2024NIMPA106969815L}. The systematic uncertainties are much smaller
than the statistical uncertainties, which are about 3\% at 2 GeV and about
7\% at 20 GeV.

{\it Results of FDs.} ---
Fig.~\ref{fig:2021FD} shows the measured CRE fluxes with a time resolution
of 6 hours in 4 energy bins for the FD event in November, 2021. Also shown
are the NM data from the OULU station (red), and $T_0$ count rate of DAMPE
(green) for the McIlwain L-values \cite{1994CASSS..10.....W} of $3.8-4.7$
which roughly corresponds to the VRC value of the OULU NM station.
The fluxes 20 days prior to the FD are used as reference for normalizations.
Note that, although the cutoff rigidities of the OULU station and the DAMPE
$T_0$ events are similar, they represent essentially different particle
energies since NMs record secondary neutrons produced by cosmic rays with
much higher energies. This explains the differences of their decrease
amplitudes.


As can be found from Fig.~\ref{fig:2021FD}, the CRE fluxes experience a fast
decrease at UT 21:00 on November 3, 2021, reach the minimum at UT 21:00 on
November 4, and gradually recover to the normal level in about 2 weeks.
The $T_0$ rate follows a trend similar to that of CRE fluxes. Different from
the September 2017 event \cite{2021ApJ...920L..43A}, no precursor was observed
for in this event. The NM profile reaches its minimum at similar time as the
$T_0$ count rate, but is about 3 hours later than the minimum time of the CRE
fluxes. The NM profile also shows a faster recovery (with characteristic
timescale of $\sim2$ days) than the CRE fluxes and the $T_0$ count rate
(with characteristic timescale of $\sim 4.5$ days).
Both the NM and the DAMPE $T_0$ data experience two-step decrease features,
as marked by arrows. Possible explanations of the two-step features could be
found in Ref.~\cite{2004JGRA..109.2117I}. We note that the $T_0$ rate exhibits
a larger decrease amplitude in the first step and a smaller amplitude in the
second step, in contrary to those of NM data. Such differences may again be
due to different energies. The time profiles for other events can be found
in Fig.~S6 of the {\tt Supplemental Material}.

The recovery phase of the FD is important in revealing the properties of the
interplanetary environment disturbed by the CME. The following function is
employed to fit the recovery profile of an FD:
\begin{equation}
\frac{I_{t}}{I_{0}}=1-A_{e} \exp\left(-\frac{t-t_{m}}{\tau}\right),
    \label{eq3}
\end{equation}
where ${I_{t}}$ is the measured CRE flux at time $t$, ${I_{0}}$ is the
normalization flux before the FD, ${A_{e}}$ is the decrease amplitude,
${t_{m}}$ is the time of minimum, and ${\tau}$ is the characteristic recovery
time.

The derived decrease amplitudes and recovery times, for each energy
bin, are shown in Fig.~\ref{fig:amp} and Fig.~\ref{fig:tau}.
Note that when performing the fitting, the systematic uncertainties
are conservatively included via a quadratic sum with the statistical
uncertainties. The decrease amplitude becomes smaller with the increasing
energy. This can be explained by the fact that particles with
higher energy are less affected by the magnetic turbulence.
The energy dependence of the recovery time, however, shows diverse
properties. For some events, the recovery time becomes shorter with the
increase of energy (e.g., the 2017-09 event), while for others it remains
nearly constant for all energies (e.g., the 2017-07 event). These diverse
behaviors may be related with properties of CMEs causing the FDs. We further
fit the energy dependence of the recovery time with a power-law function,
$\tau=aE^b$, and derive the power-law index $b$ for each event. The results
of the parameter $b$ for the 8 events are shown in Fig.~\ref{fig:tau_VOmega}.

To investigate the behaviors of the recovery time, we also show in
Fig.~\ref{fig:tau_VOmega} the possible relation between $b$ and a combined
parameter, $V_{\rm CME}\cdot\Omega$, for these events. The CME velocity
$V_{\rm CME}$ reflects the strength of the disturbance, and $\Omega$ is the
solid angle of the CME which reflects the spatial size affected by the CME.
The CME parameters are extracted from the Wang-Sheeley-Arge (WSA)-Enlil solar
wind prediction
model\footnote{https://www.swpc.noaa.gov/products/wsa-enlil-solar-wind-prediction}
\cite{1990ApJ...355..726W,2003AIPC..679..190A,2013SpWea..11...57M}.
The model reconstructs the CME speed ($V_{\text{CME}}$) and the solid angle
($\Omega$) in a region of 5 to 20 solar radii using corona graphs from three
different satellites, including STEREO-B COR2, STEREO-A COR2, and LASCO C2.
The projection effect has been taken into account. Another potentially
important parameter affecting the FDs of GCRs is the ejection direction of
the CME \cite{2016ApJ...827...13Z}. We classify the CMEs into
head-on events (I), glancing events (III), and intermediate events (II)
(see Fig.~S5 of the {\tt Supplemental Material} for a schematic plot of the
classification). Note that the boundaries among these categories are not
sharply defined, and the current classification is just for illustrative
purpose. As can be seen in Fig.~\ref{fig:tau_VOmega}, for classes I and II
events, anti-correlations between the slope $b$ and $V_{\rm CME}\cdot\Omega$
is observed. Only one class III FD event has been detected at present, and the
relationship between $b$ and $V_{\rm CME}\cdot\Omega$ is not clear yet.
An exception is the March 2022 event, which should be classified as a head-on
event based on the CME parameters (see Table~S1 of the {\tt Supplemental
Material}), but its ($b$, $V_{\rm CME}\cdot\Omega$) values place it in region
II as labelled in Fig.~\ref{fig:tau_VOmega}. The precise classification of this
event is ambiguous since the decrease amplitude is small, leading to a
relatively large uncertainty of the slope parameter $b$. Roughly, stronger CMEs
(faster in speed and wider in spread) would result in stronger energy
dependence of the recovery time.


\begin{figure}
\centering
\includegraphics[width=1.0\linewidth]{Fig4_vomega.eps}
\caption{Possible relationship between $b$ and $V_{\rm CME}\cdot\Omega$ as
revealed by the DAMPE detected FDs. The dashed lines divide the CMEs
schematically into head-on events (I), glacing events (III), and intermediate
ones (II).}
\label{fig:tau_VOmega}
\end{figure}


The transport of a charged particle in the heliosphere can be described by
the Parker's equation \cite{1965P&SS...13....9P} (see the {\tt Supplemental
Material}). CMEs and other solar activities may disturb the interplanetary
space, leading to perturbation of the propagation of GCRs. This process can
be modeled by a fast-moving diffusion barrier \cite{2018ApJ...860..160L}.
More details about the setting of the diffusion barrier can be found in
the {\tt Supplemental Material}. We adopt the stochastic differential
equation (SDE) method to solve Parker's equation with the disturbance by
CMEs \cite{1999ApJ...513..409Z,2011ApJ...735...83S}. In the modeling, the
basic parameters of the CMEs, such as the velocity, angular extension, and
direction, are from the WSA-ENLIL model. The other parameter is the diffusion
coefficient inside the barrier, which is adjusted to fit the measurements.
As shown in Fig.~S6 of the {\tt Supplemental Material}, the time
profiles of the CRE fluxes around the FDs can be reproduced by this model.
The comparison of the energy-dependence slope parameter $b$ of the recovery
time between DAMPE observations and simulations is given in Fig.~S7
of the {\tt Supplemental Material}, which also shows a good consistency.
The simulation indicates that the recovery behaviors of FDs depend on multiple
factors, including the motion and geometry of the disturbed regions. Thus the
diverse properties of the CRE FDs observed by DAMPE reflect the complicated
propagation processes of particles in the interplanetary space. Note further
that, the FD events reported in this analysis occurred between 2017 and 2024,
during which the solar magnetic polarity was always
positive\footnote{A positive solar magnetic polarity corresponds to an
outward-pointing heliospheric magnetic field in the Sun's northern hemisphere
and an inward-pointing magnetic field in the southern hemisphere.} ($A>0$).
We have a consistent drift direction configuration for all events,
isolating the additional complexity due to the polarity reversal.


{\it Conclusion.} ---
In this work we study the FDs of CREs recorded by DAMPE, from January, 2016
to March, 2024. Different from most of existing studies which focus
on the nuclear component of GCR FDs, this is the first systematic investigation
of FDs of the electron and positron component. In the total 8 FD events
detected, the largest decrease amplitude reaches about 30\%. The energy
dependence of the decrease amplitude and recovery time, from 2 GeV to 20 GeV,
is well measured. We find that the energy dependence of the recovery time show
diverse behaviors, which span from strongly energy-dependent to nearly
energy-independent. Our results further suggest that this diversity is likely
due to a combined effect of the CME velocity, angular extension, and
propagation direction. Depending on the ejection direction of the CME
relative to the Earth, an anti-correlation between the energy dependence
slope of the recovery time and the velocity times solid angle of the CME
is revealed by the data. Utilizing a diffusion barrier model, we can
account for the measurements under the framework of Parker's transport
equation, based on observed parameters of the CMEs.

With the sunspot number approaching its maximum during the solar cycle 25,
more FD events are expected to be detected, which will be very helpful in
testing the phenomenae found in this study. Particularly, additional events
belonging to with classes II and III will be especially important.
Furthermore, the comparative studies of FDs of CREs and protons or other
nuclei, as in Ref.~\cite{2018ApJ...853...76M} but with a significantly
enlarged sample are of great interest and will be carried out in future works.
Finally, the joint analysis of the FD events and the sun shadow by groundbased
cosmic ray experiments \cite{2024Innov...500695} will be crucial in improving
our understanding of impacts of solar activities on the interplanetary
magnetic environment and the particle propagation in the heliosphere.

The data that support the findings of this article are openly
available\footnote{https://doi.org/10.57760/sciencedb.space.02998}.

\section*{Acknowledgments}
The DAMPE mission was funded by the strategic priority science and technology
projects in space science of Chinese Academy of Sciences (CAS).
In China, the data analysis is supported by the National Key Research and
Development Program of China (No. 2022YFF0503302), the National Natural Science
Foundation of China (Nos. 12220101003 and 12321003), the CAS Project for Young
Scientists in Basic Research (No. YSBR-061), and the Youth Innovation Promotion
Association of CAS.
In Europe, the activities and data analysis are supported by the Swiss National
Science Foundation (SNSF), Switzerland, the National Institute for Nuclear
Physics (INFN), Italy, and the European Research Council (ERC) under the
European Union's Horizon 2020 research and innovation programme (No. 851103).

\bibliography{refs}
\bibliographystyle{apsrev}

\clearpage
\input{DAMPE_FD_SM.tex}


\end{document}
===== END FILE: DAMPE_FD.tex =====

===== BEGIN FILE: DAMPE_FD_SM.tex =====











\onecolumngrid

\setcounter{figure}{0}
\renewcommand\thefigure{S\arabic{figure}}
\setcounter{table}{0}
\renewcommand\thetable{S\arabic{table}}









\begin{center}
{\Large Supplemental Material of ``Diverse properties of electron Forbush
decreases revealed by the Dark Matter Particle Explorer''}\\
(The DAMPE collaboration)\\
\end{center}

\section{Properties of 8 CRE FDs and the associated CMEs}
Table \ref{tab:8FD} gives the main parameters of the 8 FDs discussed in this
work, including the DAMPE measured CRE fluxes, the CME parameters from the
WSA-ENLIL model, and the parameters of the diffusion barrier used in our
simulation.

\begin{table*}[!htbp]
\centering
\caption{Main parameters of the 8 FDs. From left to right, they are, the time
at which the DAMPE $T_0$ rate reaches the minimum (Time), the maximum decrease
amplitude of DAMPE CRE fluxes ($A_e^{\rm max}$), the recover time of DAMPE CRE
at 2.3 GeV ($\tau$), the slope of energy dependence of recover time ($b$), the
CME speed ($V_{\rm CME}$), the half angle width of CME ($\Theta$), the
reduction factor of the diffusion coefficient in the barrier ($\rho$),
the location of the Earth in the CME coordinate ($\theta_{0}$,$\phi_{0}$),
the latitude and longtitude angular widths of CME ($\theta_{\rm br}$,
$\phi_{\rm br}$). The symbol in parenthesis after $\phi_0$ gives the
classification of the event as defined in this work.}
\resizebox{0.8\textwidth}{22mm}{
\begin{tabular}{c|c|c|c|c|c|c|c|c|c|c}

\hline
\multicolumn{4}{c|}{DAMPE CRE Observations}&\multicolumn{2}{c}{CME Parameters}&\multicolumn{5}{|c}{Simulation Parameters} \\
\hline
\makebox[0.14\textwidth][c]{Time (UTC)} & \makebox[0.1\textwidth][c]{$A_{e}^{\rm max}$} & \makebox[0.1\textwidth][c]{$\tau$ (day)} & \makebox[0.1\textwidth][c]{$b$} & \makebox[0.12\textwidth][c]{$V_{\rm CME}$ (km s$^{-1}$)} & \makebox[0.05\textwidth][c]{$\Theta$($^{\circ}$) } & \makebox[0.04\textwidth][c]{$\rho$} & \makebox[0.04\textwidth][c]{$\theta_0$($^{\circ}$)} & \makebox[0.04\textwidth][c]{$\phi_0$($^{\circ}$)} & \makebox[0.04\textwidth][c]{$\theta_{\rm br}$($^{\circ}$)} & \makebox[0.04\textwidth][c]{$\phi_{\rm br}$($^{\circ}$)} \\
\hline
2017-07-17 04:30 & $0.125\pm0.014$ & $2.03\pm0.39$ & ~$0.010\pm0.091$ & 1199 & 57 & 4.25 & 70 & $-50$ (III) & 100 & 120 \\
2017-09-08 13:30 & $0.266\pm0.009$ & $6.46\pm0.42$ & $-0.348\pm0.054$ & 1611 & 48 & 8.25 &90 & 45 (II)  & 90  & 120 \\
2021-11-04 21:30 & $0.214\pm0.009$ & $4.69\pm0.24$ & $-0.099\pm0.065$ & 1397 & 41 & 7.00 &90 & 30 (II)  & 90  & 90 \\
2022-03-14 19:30 & $0.123\pm0.014$ & $2.61\pm0.61$ & $-0.074\pm0.018$ & 983  & 50 & 6.50 &100 & 0 (I)  & 100  & 100 \\
2023-02-28 14:30 & $0.244\pm0.011$ & $3.83\pm0.64$ & $-0.187\pm0.125$ & 1237 & 36 & 5.00 &90 & 0 (I)  & 60  & 80 \\
2023-12-01 18:30 & $0.137\pm0.027$ & $2.72\pm0.89$ & $-0.148\pm0.222$ & 944  & 43 & 5.50 &90 & 0 (I)  & 90  & 90 \\
2023-12-17 22:30 & $0.154\pm0.012$ & $4.15\pm0.56$ & $-0.355\pm0.121$ & 921  & 48 & 7.25 &90 & 10 (I)  & 100  & 90 \\
2024-03-24 22:30 & $0.254\pm0.013$ & $4.49\pm0.60$ & $-0.239\pm0.111$ & 1438 & 36 & 10.0 &100 & 0 (I)  & 90  & 80 \\
\hline
\end{tabular}
}
\label{tab:8FD}
\end{table*}

\section{Pre-selection efficiency}
The pre-selections include the trigger selection, the charge selection,
the track selection, and the shower containment requirement. The selection
efficiency as a function of energy is shown in Fig.~\ref{fig:eff}.

\begin{figure*}[!ht]
\centering
\includegraphics[width=0.6\linewidth]{FigS1_Efficiency.eps}
\caption{The CRE pre-selection effciency as a function of reconstructed energy.}
\label{fig:eff}
\end{figure*}

\section{Background subtraction}
After pre-selections, the data sample includes mainly CREs and protons.
To distinguish CREs from the proton background, a particle identification
(PID) variable is defined \cite{2021ApJ...920L..43A} to describe the shower
development in the calorimeter. Fig.~\ref{fig:PID} shows the PID distributions
in two energy bins. Curves in red and green are MC simulated PID distributions
of CREs and protons, and the blue ones are the sum of these two components,
obtained by a fit to the data using the MC templates. The proton contamination
is then obtained and subtracted from the data sample.

\begin{figure*}[!htb]
\centering
\includegraphics[width=\textwidth]{FigS2_template.eps}
\caption{The PID distributions of the flight data (black dots) and MC
simulations (red for electrons, green for protons, and blue for their sum).
The left panel is for the energy bin of $3.17 - 3.44$ GeV, and the right
panel is for the energy bin of $7.35-8.00$ GeV.}
\label{fig:PID}
\end{figure*}

\section{Exposure correction at low energies}
Due to the geomagnetic field, low energy CRs will get shielded and can not
reach the orbit of DAMPE. This corresponds to a cutoff of the low energy
spectrum of particles. To avoid the impact of the cutoff effect, a threshold
of the particle rigidity of 1.2 times of the VRC value is set, which results
in latitude-dependent (and consequently time-dependent) and energy-dependent
exposure time. For each energy and time interval, the exposure time is
calculated as the sum of the satellite travel time when this given rigidity
satisfies the above threshold condition. The time when the satellite passes
through the South Atlantic Anomaly (SAA) region, when the instrument switches
to the calibration mode, and the dead time of the data acquisation system,
is also excluded. Fig.~\ref{fig:exp} shows an illustration of the exposure time
from 2 GeV to 20 GeV for the time from November 1 to November 31, 2021. We can
see that the exposure time is almost constant above 12 GeV, and decreases at
lower energies due to the geomagnetic cutoff effect.

\begin{figure*}[!htb]
\centering
\includegraphics[width=0.7\textwidth]{FigS3_Livetime.eps}
\caption{Exposure time of DAMPE as a function of time and energy.}
\label{fig:exp}
\end{figure*}

\section{Electron transport in the heliosphere}

The transport of particles in the heliosphere is described by the Parker's
equation \cite{1965P&SS...13....9P}:
\begin{equation}
    \frac{\partial f}{\partial t} = -({\vec{V}}_{\rm sw}
    + \langle{\vec{V}}_{\rm d}\rangle)\cdot\bigtriangledown{f}
    + \bigtriangledown\cdot(K^{(\rm s)}\cdot\bigtriangledown{f})
    + \frac{1}{3}(\bigtriangledown\cdot{\vec{V}}_{\rm sw})
      \frac{\partial{f}}{\partial{\ln p}},
    \label{eq5}
\end{equation}
where $f(\vec{r},p,t)$ is the particle distribution function, $p$ is momentum,
${\vec{V}_{\rm sw}}$ is the latitude-dependent solar wind speed,
$\langle{\vec{V}}_{\rm d}\rangle$ is the pitch-angle averaged drift velocity
which has the following form \cite{1977ApJ...213..861J,1989ApJ...339..501B}
\begin{equation}
    \langle\vec{V}_{\rm d}\rangle = \frac{pv}{3q}\bigtriangledown\times\frac{\vec{B}}{B^{2}},
    \label{appeq2}
\end{equation}
where $q$ and $v$ are charge and velocity of GCR particles, and $\vec{B}$
is the interplanetary magnetic field (IMF). The variable $K^{(\rm s)}$ is the
symmetric diffusion tensor, aligned in the IMF coordinate, whose diagonal
components are the parallel diffusion coefficient $K_{||}$, and the
perpendicular diffusion coefficiens $K_{\bot\theta}$ and $K_{\bot{r}}$.

\begin{figure*}[!htb]
\centering
\includegraphics[width=0.9\textwidth]{FigS4_AMSFit.eps}
\caption{Simulation of the solar modulation effect on electrons, compared
with the AMS-02 daily electron fluxes \cite{2023PhRvL.130p1001A}, for Bartel
rotations from 2509 to 2512 in 2017.}
\label{fig:SM}
\end{figure*}

\begin{figure*}[!htb]
\centering
\includegraphics[width=0.55\linewidth]{FigS5_CME_Location.ps}
\caption{Schematic plot to show the relative location of the Earth in the CME
for the three types of events.}
\label{fig:class}
\end{figure*}

In this work the diffusion coefficients are parameterized as \cite{2019ApJ...878....6L})
\begin{equation}
  K_{i}=K_{i0}\beta\frac{B_{0}}{B}\left(\frac{P}{P_{k}}\right)^{a}
  \left[1+\left(\frac{P}{P_{k}}\right)^{\frac{b-a}{c}}\right]^c,
  \label{appeq3}
\end{equation}
where $i$ represents $||$, $\bot\theta$, or $\bot r$; $\beta$ is the particle
velocity in speed of light, $B$ is the strength of the IMF, $B_{0} \equiv 1$ nT,
$K_{i0}$ is the normalization diffusion coefficient varying with time, $p_{k}$
is the break momentum of the diffusion coefficient, $a$, $b$ are the power-law
indices before and after the break, and $c$ characterizes the smoothness of the
break.

The IMF is described by the standard Parker field as
\begin{equation}
    \vec{B}(r,\theta) = \frac{AB_{\oplus}}{r^{2}}\left(\hat{e}_{r}-\frac{r\Omega{\sin\theta}}{V_{\rm sw}}\hat{e}_{\theta}\right) [1-2H(\theta-\theta_{\rm cs})],
    \label{addeq5}
\end{equation}
where $B_{\oplus}$ is an average of 13 months magnetic field strength before
the specific time at 1 AU, $A$ describes the polarity of the IMF, defined as
$A=\pm{1}/\sqrt{1+(\Omega/V_{\rm sw})^{2}}$, $\theta_{\rm cs}$ determines the
polar extent of the Heliospheric Current Sheet (HCS), and $\Omega$ is the
angular velocity of the sun.

The Parker's equation is solved numerically with the stochastic differential
equation (SDE) technique \cite{1999ApJ...513..409Z,2011ApJ...735...83S}. The
parameters of the diffusion coefficient are derived through fitting the
electron fluxes measured by AMS-02 \cite{2023PhRvL.130p1001A}, using the local
interstellar spectrum of electrons derived in Ref.~\cite{2021APh...12402495Z}
(see Fig.~\ref{fig:SM} for an example of the model fitting results in 2017).
The relatively small portion of positrons has been omitted in the calculation.

\section{The diffusion barrier model of FD}
The CME is described by a three dimensional diffusion barrier, in which the
diffusion coefficient is smaller than the interplanetary value
\cite{2018ApJ...860..160L}. The diffusion coefficient inside the barrier is
written as
\begin{equation}
    K{'}_{i}=\frac{K_{i}}{1+\rho{h(\theta)f(r)g(\phi)}},
\end{equation}
where $\rho$ is the reduction factor, $h(\theta), g(\phi)$, and
$f(r)$ describes the geometry of the diffusion barrier, as
\begin{equation}
    h(\theta)=e^{-\left(\frac{\theta-\theta_{0}}{\theta_{\rm br}}\right)^{10}},
\end{equation}
\begin{equation}
    g(\phi)=e^{-\left(\frac{\phi-\phi_{0}}{\phi_{\rm br}}\right)^{10}},
\end{equation}
\begin{equation}
 f(r)=\left\{
 \begin{aligned}
 1-\frac{r-r_{\rm cen}}{r_{a}}&,&r_{\rm cen}<r<r_{\rm fr} \\
 \frac{r-r_{\rm end}}{r_b}&,& r_{\rm end}<r\leq{r_{\rm cen}}\\
 0& ,& {\rm others}
 \end{aligned}
 \right.,
\end{equation}
where $(\theta_0,\phi_0)$ are the direction of the CME center.
For $\theta_0=\pi/2$ and $\phi_0=0$, the CME hits the Earth head on.
Parameter $\theta_{\rm br}$ and $\phi_{\rm br}$ describes the extension
of the barrier. The radical location of the front, center, and end of the
CME are depicted as $r_{\rm fr}$, $r_{\rm cen}$, and $r_{\rm end}$, and
the widths of the leading and trailing parts of the barrier are
$r_{a}=r_{\rm fr}-r_{\rm cen}$ and $r_{b}=r_{\rm cen}-r_{\rm end}$.
The same SDE method is employed to calculate the time-dependent electron
fluxes including the diffusion barrier. The parameters used to fit the
DAMPE data are given in Table \ref{tab:8FD}. A comparison between the
observed time profiles and simulated ones for the 8 FDs is shown in
Fig.~\ref{fig:simu}. We also fit the slopes of the energy-dependence of
the recovery time for the simulation results, which are shown in
Fig.~\ref{fig:b} with a comparison with observations. Here the
uncertainties of $b_{\rm Simu}$ are obtained through scaling down
the simulation sample to similar event statistics with the flight
data, which makes the comparison between the observations and
simulations more direct.


\begin{figure*}[!htb]
\centering
\includegraphics[width=0.9\textwidth]{FigS6_SimuA.eps}
\includegraphics[width=0.9\textwidth]{FigS6_SimuB.eps}
\caption{Comparison of the DAMPE observed CRE time profiles (dots) with the
simulation results (lines) for 5 selected energy bins.}
\label{fig:simu}
\end{figure*}

\begin{figure*}[!htb]
\centering
\includegraphics[width=0.7\linewidth]{FigS7_parab.eps}
\caption{The slope parameter $b$ derived from DAMPE observations
($b_{\rm Obs}$) and simulations ($b_{\rm Simu}$).}
\label{fig:b}
\end{figure*}
===== END FILE: DAMPE_FD_SM.tex =====

===== BEGIN FILE: author20250812.tex =====
\author{F.~Alemanno}
\affiliation{Dipartimento di Matematica e Fisica E. De Giorgi, Universit\`a del Salento, I-73100, Lecce, Italy}
\affiliation{Istituto Nazionale di Fisica Nucleare (INFN) - Sezione di Lecce, I-73100, Lecce, Italy}

\author{Q.~An}
\affiliation{State Key Laboratory of Particle Detection and Electronics, University of Science and Technology of China, Hefei 230026, China}
\affiliation{Department of Modern Physics, University of Science and Technology of China, Hefei 230026, China}

\author{P.~Azzarello}
\affiliation{Department of Nuclear and Particle Physics, University of Geneva, CH-1211, Switzerland}

\author{F.~C.~T.~Barbato}
\affiliation{Gran Sasso Science Institute (GSSI), Via Iacobucci 2, I-67100 L’Aquila, Italy}
\affiliation{Istituto Nazionale di Fisica Nucleare (INFN) - Laboratori Nazionali del Gran Sasso, I-67100 Assergi, L’Aquila, Italy}

\author{P.~Bernardini}
\affiliation{Dipartimento di Matematica e Fisica E. De Giorgi, Universit\`a del Salento, I-73100, Lecce, Italy}
\affiliation{Istituto Nazionale di Fisica Nucleare (INFN) - Sezione di Lecce, I-73100, Lecce, Italy}

\author{X.~J.~Bi}
\affiliation{University of Chinese Academy of Sciences, Beijing 100049, China}
\affiliation{Particle Astrophysics Division, Institute of High Energy Physics, Chinese Academy of Sciences, Beijing 100049, China}

\author{H.~Boutin}
\affiliation{Department of Nuclear and Particle Physics, University of Geneva, CH-1211, Switzerland}

\author{I.~Cagnoli}
\affiliation{Gran Sasso Science Institute (GSSI), Via Iacobucci 2, I-67100 L’Aquila, Italy}
\affiliation{Istituto Nazionale di Fisica Nucleare (INFN) - Laboratori Nazionali del Gran Sasso, I-67100 Assergi, L’Aquila, Italy}

\author{M.~S.~Cai}
\affiliation{Key Laboratory of Dark Matter and Space Astronomy, Purple Mountain Observatory, Chinese Academy of Sciences, Nanjing 210023, China}
\affiliation{School of Astronomy and Space Science, University of Science and Technology of China, Hefei 230026, China}

\author{E.~Casilli}
\affiliation{Gran Sasso Science Institute (GSSI), Via Iacobucci 2, I-67100 L’Aquila, Italy}
\affiliation{Istituto Nazionale di Fisica Nucleare (INFN) - Laboratori Nazionali del Gran Sasso, I-67100 Assergi, L’Aquila, Italy}

\author{J.~Chang}
\affiliation{Key Laboratory of Dark Matter and Space Astronomy, Purple Mountain Observatory, Chinese Academy of Sciences, Nanjing 210023, China}
\affiliation{School of Astronomy and Space Science, University of Science and Technology of China, Hefei 230026, China}

\author{D.~Y.~Chen}
\affiliation{Key Laboratory of Dark Matter and Space Astronomy, Purple Mountain Observatory, Chinese Academy of Sciences, Nanjing 210023, China}

\author{J.~L~Chen}
\affiliation{Institute of Modern Physics, Chinese Academy of Sciences, Lanzhou 730000, China}

\author{Z.~F.~Chen}
\affiliation{Institute of Modern Physics, Chinese Academy of Sciences, Lanzhou 730000, China}

\author{Z.X.~Chen}
\affiliation{Institute of Modern Physics, Chinese Academy of Sciences, Lanzhou 730000, China}
\affiliation{University of Chinese Academy of Sciences, Beijing 100049, China}

\author{P.~Coppin}
\affiliation{Department of Nuclear and Particle Physics, University of Geneva, CH-1211, Switzerland}

\author{M.~Y.~Cui}
\affiliation{Key Laboratory of Dark Matter and Space Astronomy, Purple Mountain Observatory, Chinese Academy of Sciences, Nanjing 210023, China}

\author{T.~S.~Cui}
\affiliation{National Space Science Center, Chinese Academy of Sciences, Beijing 100190, China}

\author{I.~De Mitri}
\affiliation{Gran Sasso Science Institute (GSSI), Via Iacobucci 2, I-67100 L’Aquila, Italy}
\affiliation{Istituto Nazionale di Fisica Nucleare (INFN) - Laboratori Nazionali del Gran Sasso, I-67100 Assergi, L’Aquila, Italy}

\author{F.~de Palma}
\affiliation{Dipartimento di Matematica e Fisica E. De Giorgi, Universit\`a del Salento, I-73100, Lecce, Italy}
\affiliation{Istituto Nazionale di Fisica Nucleare (INFN) - Sezione di Lecce, I-73100, Lecce, Italy}

\author{A.~Di Giovanni}
\affiliation{Gran Sasso Science Institute (GSSI), Via Iacobucci 2, I-67100 L’Aquila, Italy}
\affiliation{Istituto Nazionale di Fisica Nucleare (INFN) - Laboratori Nazionali del Gran Sasso, I-67100 Assergi, L’Aquila, Italy}

\author{T.~K.~Dong}
\affiliation{Key Laboratory of Dark Matter and Space Astronomy, Purple Mountain Observatory, Chinese Academy of Sciences, Nanjing 210023, China}

\author{Z.~X.~Dong}
\affiliation{National Space Science Center, Chinese Academy of Sciences, Beijing 100190, China}

\author{G.~Donvito}
\affiliation{Istituto Nazionale di Fisica Nucleare, Sezione di Bari, via Orabona 4, I-70126 Bari, Italy}

\author{J.~L.~Duan}
\affiliation{Institute of Modern Physics, Chinese Academy of Sciences, Lanzhou 730000, China}

\author{K.~K.~Duan}
\affiliation{Key Laboratory of Dark Matter and Space Astronomy, Purple Mountain Observatory, Chinese Academy of Sciences, Nanjing 210023, China}

\author{R.~R.~Fan}
\affiliation{Particle Astrophysics Division, Institute of High Energy Physics, Chinese Academy of Sciences, Beijing 100049, China}

\author{Y.~Z.~Fan}
\affiliation{Key Laboratory of Dark Matter and Space Astronomy, Purple Mountain Observatory, Chinese Academy of Sciences, Nanjing 210023, China}
\affiliation{School of Astronomy and Space Science, University of Science and Technology of China, Hefei 230026, China}

\author{F.~Fang}
\affiliation{Institute of Modern Physics, Chinese Academy of Sciences, Lanzhou 730000, China}

\author{K.~Fang}
\affiliation{Particle Astrophysics Division, Institute of High Energy Physics, Chinese Academy of Sciences, Beijing 100049, China}

\author{C.~Q.~Feng}
\affiliation{State Key Laboratory of Particle Detection and Electronics, University of Science and Technology of China, Hefei 230026, China}
\affiliation{Department of Modern Physics, University of Science and Technology of China, Hefei 230026, China}

\author{L.~Feng}
\affiliation{Key Laboratory of Dark Matter and Space Astronomy, Purple Mountain Observatory, Chinese Academy of Sciences, Nanjing 210023, China}

\author{S.~Fogliacco}
\affiliation{Gran Sasso Science Institute (GSSI), Via Iacobucci 2, I-67100 L’Aquila, Italy}
\affiliation{Istituto Nazionale di Fisica Nucleare (INFN) - Laboratori Nazionali del Gran Sasso, I-67100 Assergi, L’Aquila, Italy}

\author{J.~M.~Frieden}
\altaffiliation{Now at Institute of Physics, Ecole Polytechnique F\'{e}d\'{e}rale de Lausanne (EPFL), CH-1015 Lausanne, Switzerland.}
\affiliation{Department of Nuclear and Particle Physics, University of Geneva, CH-1211, Switzerland}

\author{P.~Fusco}
\affiliation{Istituto Nazionale di Fisica Nucleare, Sezione di Bari, via Orabona 4, I-70126 Bari, Italy}
\affiliation{Dipartimento di Fisica ``M.~Merlin'', dell’Universit\`a e del Politecnico di Bari, via Amendola 173, I-70126 Bari, Italy}

\author{M.~Gao}
\affiliation{Particle Astrophysics Division, Institute of High Energy Physics, Chinese Academy of Sciences, Beijing 100049, China}

\author{F.~Gargano}
\affiliation{Istituto Nazionale di Fisica Nucleare, Sezione di Bari, via Orabona 4, I-70126 Bari, Italy}

\author{E.~Ghose}
\affiliation{Dipartimento di Matematica e Fisica E. De Giorgi, Universit\`a del Salento, I-73100, Lecce, Italy}
\affiliation{Istituto Nazionale di Fisica Nucleare (INFN) - Sezione di Lecce, I-73100, Lecce, Italy}

\author{K.~Gong}
\affiliation{Particle Astrophysics Division, Institute of High Energy Physics, Chinese Academy of Sciences, Beijing 100049, China}

\author{Y.~Z.~Gong}
\affiliation{Key Laboratory of Dark Matter and Space Astronomy, Purple Mountain Observatory, Chinese Academy of Sciences, Nanjing 210023, China}

\author{D.~Y.~Guo}
\affiliation{Particle Astrophysics Division, Institute of High Energy Physics, Chinese Academy of Sciences, Beijing 100049, China}

\author{J.~H.~Guo}
\affiliation{Key Laboratory of Dark Matter and Space Astronomy, Purple Mountain Observatory, Chinese Academy of Sciences, Nanjing 210023, China}
\affiliation{School of Astronomy and Space Science, University of Science and Technology of China, Hefei 230026, China}

\author{S.~X.~Han}
\affiliation{National Space Science Center, Chinese Academy of Sciences, Beijing 100190, China}

\author{Y.~M.~Hu}
\affiliation{Key Laboratory of Dark Matter and Space Astronomy, Purple Mountain Observatory, Chinese Academy of Sciences, Nanjing 210023, China}

\author{G.~S.~Huang}
\affiliation{State Key Laboratory of Particle Detection and Electronics, University of Science and Technology of China, Hefei 230026, China}
\affiliation{Department of Modern Physics, University of Science and Technology of China, Hefei 230026, China}

\author{X.~Y.~Huang}
\affiliation{Key Laboratory of Dark Matter and Space Astronomy, Purple Mountain Observatory, Chinese Academy of Sciences, Nanjing 210023, China}
\affiliation{School of Astronomy and Space Science, University of Science and Technology of China, Hefei 230026, China}

\author{Y.~Y.~Huang}
\affiliation{Key Laboratory of Dark Matter and Space Astronomy, Purple Mountain Observatory, Chinese Academy of Sciences, Nanjing 210023, China}

\author{M.~Ionica}
\affiliation{Istituto Nazionale di Fisica Nucleare (INFN) - Sezione di Perugia, I-06123 Perugia, Italy}

\author{L.~Y.~Jiang}
\affiliation{Key Laboratory of Dark Matter and Space Astronomy, Purple Mountain Observatory, Chinese Academy of Sciences, Nanjing 210023, China}

\author{W.~Jiang}
\affiliation{Key Laboratory of Dark Matter and Space Astronomy, Purple Mountain Observatory, Chinese Academy of Sciences, Nanjing 210023, China}

\author{Y.Z.~Jiang}
\altaffiliation{Also at Dipartimento di Fisica e Geologia, Universit\`a degli Studi di Perugia, I-06123 Perugia, Italy.}
\affiliation{Istituto Nazionale di Fisica Nucleare (INFN) - Sezione di Perugia, I-06123 Perugia, Italy}

\author{J.~Kong}
\affiliation{Institute of Modern Physics, Chinese Academy of Sciences, Lanzhou 730000, China}

\author{A.~Kotenko}
\affiliation{Department of Nuclear and Particle Physics, University of Geneva, CH-1211, Switzerland}

\author{D.~Kyratzis}
\affiliation{Gran Sasso Science Institute (GSSI), Via Iacobucci 2, I-67100 L’Aquila, Italy}
\affiliation{Istituto Nazionale di Fisica Nucleare (INFN) - Laboratori Nazionali del Gran Sasso, I-67100 Assergi, L’Aquila, Italy}

\author{S.~J.~Lei}
\affiliation{Key Laboratory of Dark Matter and Space Astronomy, Purple Mountain Observatory, Chinese Academy of Sciences, Nanjing 210023, China}

\author{B.~Li}
\affiliation{Key Laboratory of Dark Matter and Space Astronomy, Purple Mountain Observatory, Chinese Academy of Sciences, Nanjing 210023, China}
\affiliation{School of Astronomy and Space Science, University of Science and Technology of China, Hefei 230026, China}

\author{W.~L.~Li}
\affiliation{National Space Science Center, Chinese Academy of Sciences, Beijing 100190, China}

\author{W.~H.~Li}
\affiliation{Key Laboratory of Dark Matter and Space Astronomy, Purple Mountain Observatory, Chinese Academy of Sciences, Nanjing 210023, China}

\author{X.~Li}
\affiliation{Key Laboratory of Dark Matter and Space Astronomy, Purple Mountain Observatory, Chinese Academy of Sciences, Nanjing 210023, China}
\affiliation{School of Astronomy and Space Science, University of Science and Technology of China, Hefei 230026, China}

\author{X.~Q.~Li}
\affiliation{National Space Science Center, Chinese Academy of Sciences, Beijing 100190, China}

\author{Y.~M.~Liang}
\affiliation{National Space Science Center, Chinese Academy of Sciences, Beijing 100190, China}

\author{C.M.~Liu}
\affiliation{Istituto Nazionale di Fisica Nucleare (INFN) - Sezione di Perugia, I-06123 Perugia, Italy}

\author{H.~Liu}
\affiliation{Key Laboratory of Dark Matter and Space Astronomy, Purple Mountain Observatory, Chinese Academy of Sciences, Nanjing 210023, China}

\author{J.~Liu}
\affiliation{Institute of Modern Physics, Chinese Academy of Sciences, Lanzhou 730000, China}

\author{S.~B.~Liu}
\affiliation{State Key Laboratory of Particle Detection and Electronics, University of Science and Technology of China, Hefei 230026, China}
\affiliation{Department of Modern Physics, University of Science and Technology of China, Hefei 230026, China}

\author{Y.~Liu}
\affiliation{Key Laboratory of Dark Matter and Space Astronomy, Purple Mountain Observatory, Chinese Academy of Sciences, Nanjing 210023, China}

\author{F.~Loparco}
\affiliation{Istituto Nazionale di Fisica Nucleare, Sezione di Bari, via Orabona 4, I-70126 Bari, Italy}
\affiliation{Dipartimento di Fisica ``M.~Merlin'', dell’Universit\`a e del Politecnico di Bari, via Amendola 173, I-70126 Bari, Italy}

\author{M.~Ma}
\affiliation{National Space Science Center, Chinese Academy of Sciences, Beijing 100190, China}

\author{P.~X.~Ma}
\affiliation{Key Laboratory of Dark Matter and Space Astronomy, Purple Mountain Observatory, Chinese Academy of Sciences, Nanjing 210023, China}

\author{T.~Ma}
\affiliation{Key Laboratory of Dark Matter and Space Astronomy, Purple Mountain Observatory, Chinese Academy of Sciences, Nanjing 210023, China}

\author{X.~Y.~Ma}
\affiliation{National Space Science Center, Chinese Academy of Sciences, Beijing 100190, China}

\author{G.~Marsella}
\altaffiliation{Now at Dipartimento di Fisica e Chimica ``E. Segr\`e'', Universit\`a degli Studi di Palermo, via delle Scienze ed. 17, I-90128 Palermo, Italy.}
\affiliation{Dipartimento di Matematica e Fisica E. De Giorgi, Universit\`a del Salento, I-73100, Lecce, Italy}
\affiliation{Istituto Nazionale di Fisica Nucleare (INFN) - Sezione di Lecce, I-73100, Lecce, Italy}

\author{M.~N.~Mazziotta}
\affiliation{Istituto Nazionale di Fisica Nucleare, Sezione di Bari, via Orabona 4, I-70126 Bari, Italy}

\author{D.~Mo}
\affiliation{Institute of Modern Physics, Chinese Academy of Sciences, Lanzhou 730000, China}

\author{Y.~Nie}
\affiliation{State Key Laboratory of Particle Detection and Electronics, University of Science and Technology of China, Hefei 230026, China}
\affiliation{Department of Modern Physics, University of Science and Technology of China, Hefei 230026, China}

\author{X.~Y.~Niu}
\affiliation{Institute of Modern Physics, Chinese Academy of Sciences, Lanzhou 730000, China}

\author{A.~Parenti}
\altaffiliation{Now at Inter-university Institute for High Energies, Universit\`e Libre de Bruxelles, B-1050 Brussels, Belgium.}
\affiliation{Gran Sasso Science Institute (GSSI), Via Iacobucci 2, I-67100 L’Aquila, Italy}
\affiliation{Istituto Nazionale di Fisica Nucleare (INFN) - Laboratori Nazionali del Gran Sasso, I-67100 Assergi, L’Aquila, Italy}

\author{W.~X.~Peng}
\affiliation{Particle Astrophysics Division, Institute of High Energy Physics, Chinese Academy of Sciences, Beijing 100049, China}

\author{X.~Y.~Peng}
\affiliation{Key Laboratory of Dark Matter and Space Astronomy, Purple Mountain Observatory, Chinese Academy of Sciences, Nanjing 210023, China}

\author{C.~Perrina}
\altaffiliation{Now at Institute of Physics, Ecole Polytechnique F\'{e}d\'{e}rale de Lausanne (EPFL), CH-1015 Lausanne, Switzerland.}
\affiliation{Department of Nuclear and Particle Physics, University of Geneva, CH-1211, Switzerland}

\author{E.~Putti-Garcia}
\affiliation{Department of Nuclear and Particle Physics, University of Geneva, CH-1211, Switzerland}

\author{R.~Qiao}
\affiliation{Particle Astrophysics Division, Institute of High Energy Physics, Chinese Academy of Sciences, Beijing 100049, China}

\author{J.~N.~Rao}
\affiliation{National Space Science Center, Chinese Academy of Sciences, Beijing 100190, China}

\author{Y.~Rong}
\affiliation{State Key Laboratory of Particle Detection and Electronics, University of Science and Technology of China, Hefei 230026, China}
\affiliation{Department of Modern Physics, University of Science and Technology of China, Hefei 230026, China}

\author{A.~Ruina}
\affiliation{Department of Nuclear and Particle Physics, University of Geneva, CH-1211, Switzerland}

\author{R.~Sarkar}
\affiliation{Gran Sasso Science Institute (GSSI), Via Iacobucci 2, I-67100 L’Aquila, Italy}
\affiliation{Istituto Nazionale di Fisica Nucleare (INFN) - Laboratori Nazionali del Gran Sasso, I-67100 Assergi, L’Aquila, Italy}

\author{P.~Savina}
\affiliation{Gran Sasso Science Institute (GSSI), Via Iacobucci 2, I-67100 L’Aquila, Italy}
\affiliation{Istituto Nazionale di Fisica Nucleare (INFN) - Laboratori Nazionali del Gran Sasso, I-67100 Assergi, L’Aquila, Italy}

\author{Z.~Shangguan}
\affiliation{National Space Science Center, Chinese Academy of Sciences, Beijing 100190, China}

\author{W.~H.~Shen}
\affiliation{National Space Science Center, Chinese Academy of Sciences, Beijing 100190, China}

\author{Z.~Q.~Shen}
\affiliation{Key Laboratory of Dark Matter and Space Astronomy, Purple Mountain Observatory, Chinese Academy of Sciences, Nanjing 210023, China}

\author{Z.~T.~Shen}
\affiliation{State Key Laboratory of Particle Detection and Electronics, University of Science and Technology of China, Hefei 230026, China}
\affiliation{Department of Modern Physics, University of Science and Technology of China, Hefei 230026, China}

\author{L.~Silveri}
\altaffiliation{Now at New York University Abu Dhabi, Saadiyat Island, Abu Dhabi 129188, United Arab Emirates.}
\affiliation{Gran Sasso Science Institute (GSSI), Via Iacobucci 2, I-67100 L’Aquila, Italy}
\affiliation{Istituto Nazionale di Fisica Nucleare (INFN) - Laboratori Nazionali del Gran Sasso, I-67100 Assergi, L’Aquila, Italy}

\author{J.~X.~Song}
\affiliation{National Space Science Center, Chinese Academy of Sciences, Beijing 100190, China}

\author{M.~Stolpovskiy}
\affiliation{Department of Nuclear and Particle Physics, University of Geneva, CH-1211, Switzerland}

\author{H.~Su}
\affiliation{Institute of Modern Physics, Chinese Academy of Sciences, Lanzhou 730000, China}

\author{M.~Su}
\affiliation{Department of Physics and Laboratory for Space Research, University of Hong Kong, Hong Kong SAR 999077, China}

\author{H.~R.~Sun}
\affiliation{State Key Laboratory of Particle Detection and Electronics, University of Science and Technology of China, Hefei 230026, China}
\affiliation{Department of Modern Physics, University of Science and Technology of China, Hefei 230026, China}

\author{Z.~Y.~Sun}
\affiliation{Institute of Modern Physics, Chinese Academy of Sciences, Lanzhou 730000, China}

\author{A.~Surdo}
\affiliation{Istituto Nazionale di Fisica Nucleare (INFN) - Sezione di Lecce, I-73100, Lecce, Italy}

\author{X.~J.~Teng}
\affiliation{National Space Science Center, Chinese Academy of Sciences, Beijing 100190, China}

\author{A.~Tykhonov}
\affiliation{Department of Nuclear and Particle Physics, University of Geneva, CH-1211, Switzerland}

\author{G.~F.~Wang}
\affiliation{State Key Laboratory of Particle Detection and Electronics, University of Science and Technology of China, Hefei 230026, China}
\affiliation{Department of Modern Physics, University of Science and Technology of China, Hefei 230026, China}

\author{J.~Z.~Wang}
\affiliation{Particle Astrophysics Division, Institute of High Energy Physics, Chinese Academy of Sciences, Beijing 100049, China}

\author{L.~G.~Wang}
\affiliation{National Space Science Center, Chinese Academy of Sciences, Beijing 100190, China}

\author{S.~Wang}
\affiliation{Key Laboratory of Dark Matter and Space Astronomy, Purple Mountain Observatory, Chinese Academy of Sciences, Nanjing 210023, China}

\author{X.~L.~Wang}
\affiliation{State Key Laboratory of Particle Detection and Electronics, University of Science and Technology of China, Hefei 230026, China}
\affiliation{Department of Modern Physics, University of Science and Technology of China, Hefei 230026, China}

\author{Y.~F.~Wang}
\affiliation{State Key Laboratory of Particle Detection and Electronics, University of Science and Technology of China, Hefei 230026, China}
\affiliation{Department of Modern Physics, University of Science and Technology of China, Hefei 230026, China}

\author{D.~M.~Wei}
\affiliation{Key Laboratory of Dark Matter and Space Astronomy, Purple Mountain Observatory, Chinese Academy of Sciences, Nanjing 210023, China}
\affiliation{School of Astronomy and Space Science, University of Science and Technology of China, Hefei 230026, China}

\author{J.~J.~Wei}
\affiliation{Key Laboratory of Dark Matter and Space Astronomy, Purple Mountain Observatory, Chinese Academy of Sciences, Nanjing 210023, China}

\author{Y.~F.~Wei}
\affiliation{State Key Laboratory of Particle Detection and Electronics, University of Science and Technology of China, Hefei 230026, China}
\affiliation{Department of Modern Physics, University of Science and Technology of China, Hefei 230026, China}

\author{D.~Wu}
\affiliation{Particle Astrophysics Division, Institute of High Energy Physics, Chinese Academy of Sciences, Beijing 100049, China}

\author{J.~Wu}
\affiliation{Key Laboratory of Dark Matter and Space Astronomy, Purple Mountain Observatory, Chinese Academy of Sciences, Nanjing 210023, China}
\affiliation{School of Astronomy and Space Science, University of Science and Technology of China, Hefei 230026, China}

\author{S.~S.~Wu}
\affiliation{National Space Science Center, Chinese Academy of Sciences, Beijing 100190, China}

\author{X.~Wu}
\affiliation{Department of Nuclear and Particle Physics, University of Geneva, CH-1211, Switzerland}

\author{Z.~Q.~Xia}
\affiliation{Key Laboratory of Dark Matter and Space Astronomy, Purple Mountain Observatory, Chinese Academy of Sciences, Nanjing 210023, China}

\author{Z.~Xiong}
\affiliation{Gran Sasso Science Institute (GSSI), Via Iacobucci 2, I-67100 L’Aquila, Italy}
\affiliation{Istituto Nazionale di Fisica Nucleare (INFN) - Laboratori Nazionali del Gran Sasso, I-67100 Assergi, L’Aquila, Italy}

\author{E.~H.~Xu}
\affiliation{State Key Laboratory of Particle Detection and Electronics, University of Science and Technology of China, Hefei 230026, China}
\affiliation{Department of Modern Physics, University of Science and Technology of China, Hefei 230026, China}

\author{H.~T.~Xu}
\affiliation{National Space Science Center, Chinese Academy of Sciences, Beijing 100190, China}

\author{J.~Xu}
\affiliation{Key Laboratory of Dark Matter and Space Astronomy, Purple Mountain Observatory, Chinese Academy of Sciences, Nanjing 210023, China}

\author{Z.~H.~Xu}
\affiliation{Institute of Modern Physics, Chinese Academy of Sciences, Lanzhou 730000, China}

\author{Z.~Z.~Xu}
\affiliation{State Key Laboratory of Particle Detection and Electronics, University of Science and Technology of China, Hefei 230026, China}
\affiliation{Department of Modern Physics, University of Science and Technology of China, Hefei 230026, China}

\author{Z.~L.~Xu}
\affiliation{Key Laboratory of Dark Matter and Space Astronomy, Purple Mountain Observatory, Chinese Academy of Sciences, Nanjing 210023, China}

\author{G.~F.~Xue}
\affiliation{National Space Science Center, Chinese Academy of Sciences, Beijing 100190, China}

\author{M.~Y.~Yan}
\affiliation{State Key Laboratory of Particle Detection and Electronics, University of Science and Technology of China, Hefei 230026, China}
\affiliation{Department of Modern Physics, University of Science and Technology of China, Hefei 230026, China}

\author{H.~B.~Yang}
\affiliation{Institute of Modern Physics, Chinese Academy of Sciences, Lanzhou 730000, China}

\author{P.~Yang}
\affiliation{Institute of Modern Physics, Chinese Academy of Sciences, Lanzhou 730000, China}

\author{Y.~Q.~Yang}
\affiliation{Institute of Modern Physics, Chinese Academy of Sciences, Lanzhou 730000, China}

\author{H.~J.~Yao}
\affiliation{Institute of Modern Physics, Chinese Academy of Sciences, Lanzhou 730000, China}

\author{Y.~H.~Yu}
\affiliation{Institute of Modern Physics, Chinese Academy of Sciences, Lanzhou 730000, China}

\author{Q.~Yuan}
\affiliation{Key Laboratory of Dark Matter and Space Astronomy, Purple Mountain Observatory, Chinese Academy of Sciences, Nanjing 210023, China}
\affiliation{School of Astronomy and Space Science, University of Science and Technology of China, Hefei 230026, China}

\author{C.~Yue}
\affiliation{Key Laboratory of Dark Matter and Space Astronomy, Purple Mountain Observatory, Chinese Academy of Sciences, Nanjing 210023, China}

\author{J.~J.~Zang}
\altaffiliation{Also at School of Physics and Electronic Engineering, Linyi University, Linyi 276000, China.}
\affiliation{Key Laboratory of Dark Matter and Space Astronomy, Purple Mountain Observatory, Chinese Academy of Sciences, Nanjing 210023, China}

\author{S.~X.~Zhang}
\affiliation{Institute of Modern Physics, Chinese Academy of Sciences, Lanzhou 730000, China}

\author{W.~Z.~Zhang}
\affiliation{National Space Science Center, Chinese Academy of Sciences, Beijing 100190, China}

\author{Yan~Zhang}
\affiliation{Key Laboratory of Dark Matter and Space Astronomy, Purple Mountain Observatory, Chinese Academy of Sciences, Nanjing 210023, China}

\author{Y.~P.~Zhang}
\affiliation{Institute of Modern Physics, Chinese Academy of Sciences, Lanzhou 730000, China}

\author{Yi~Zhang}
\affiliation{Key Laboratory of Dark Matter and Space Astronomy, Purple Mountain Observatory, Chinese Academy of Sciences, Nanjing 210023, China}
\affiliation{School of Astronomy and Space Science, University of Science and Technology of China, Hefei 230026, China}

\author{Y.~J.~Zhang}
\affiliation{Institute of Modern Physics, Chinese Academy of Sciences, Lanzhou 730000, China}

\author{Y.~Q.~Zhang}
\affiliation{Key Laboratory of Dark Matter and Space Astronomy, Purple Mountain Observatory, Chinese Academy of Sciences, Nanjing 210023, China}

\author{Y.~L.~Zhang}
\affiliation{State Key Laboratory of Particle Detection and Electronics, University of Science and Technology of China, Hefei 230026, China}
\affiliation{Department of Modern Physics, University of Science and Technology of China, Hefei 230026, China}

\author{Z.~Zhang}
\affiliation{Key Laboratory of Dark Matter and Space Astronomy, Purple Mountain Observatory, Chinese Academy of Sciences, Nanjing 210023, China}

\author{Z.~Y.~Zhang}
\affiliation{State Key Laboratory of Particle Detection and Electronics, University of Science and Technology of China, Hefei 230026, China}
\affiliation{Department of Modern Physics, University of Science and Technology of China, Hefei 230026, China}

\author{C.~Zhao}
\affiliation{State Key Laboratory of Particle Detection and Electronics, University of Science and Technology of China, Hefei 230026, China}
\affiliation{Department of Modern Physics, University of Science and Technology of China, Hefei 230026, China}

\author{H.~Y.~Zhao}
\affiliation{Institute of Modern Physics, Chinese Academy of Sciences, Lanzhou 730000, China}

\author{X.~F.~Zhao}
\affiliation{National Space Science Center, Chinese Academy of Sciences, Beijing 100190, China}

\author{C.~Y.~Zhou}
\affiliation{National Space Science Center, Chinese Academy of Sciences, Beijing 100190, China}

\author{X.~Zhu}
\altaffiliation{Also at School of computing, Nanjing University of Posts and Telecommunications, Nanjing 210023, China.}
\affiliation{Key Laboratory of Dark Matter and Space Astronomy, Purple Mountain Observatory, Chinese Academy of Sciences, Nanjing 210023, China}

\author{Y.~Zhu}
\affiliation{National Space Science Center, Chinese Academy of Sciences, Beijing 100190, China}


\collaboration{DAMPE Collaboration}\email{dampe@pmo.ac.cn}

\author{X.~Luo}
\affiliation{Shandong Institute of Advanced Technology, Jinan 250100, China}
===== END FILE: author20250812.tex =====
